% ==STYLE-BEGIN==
\documentclass[UTF8,11pt]{ctexart}
\usepackage{amsmath,amssymb,amsthm}
\usepackage[a4paper,margin=2.4cm,top=2.3cm,bottom=2.3cm]{geometry}
\usepackage{fancyhdr}
\usepackage{enumitem}
\usepackage{hyperref}
\hypersetup{hidelinks}

% ---- 版式 ----
\linespread{1.15}
\setlength{\parskip}{0.35em}
\setlength{\parindent}{1.6em}
\setlist{nosep,leftmargin=2em}

% ---- 页眉 ----
\pagestyle{fancy}
\fancyhf{}
\fancyhead[L]{\small\coursename}
\fancyhead[R]{\small\thepage}
\renewcommand{\headrulewidth}{0.4pt}

% ---- 定理环境：中文用正体，标题加粗，按小节共用一条流水号 ----
\newtheoremstyle{cnthm}{0.7em}{0.7em}{\normalfont}{}{\bfseries}{}{0.5em}{}
\theoremstyle{cnthm}
\newtheorem{thm}{定理}[section]
\newtheorem{lem}[thm]{引理}
\newtheorem{prop}[thm]{命题}
\newtheorem{cor}[thm]{推论}
\newtheorem{defn}[thm]{定义}
\newtheorem{rem}[thm]{注}
\newtheorem{exa}[thm]{例}
\newtheorem{conj}[thm]{猜想}
\renewcommand{\proofname}{\bfseries 证明}
\providecommand{\coursename}{}
% ==STYLE-END==
\renewcommand{\coursename}{高等代数\,·\,1111}
\xeCJKDeclareCharClass{CJK}{"2150 -> "218F}
\title{高等代数 1111}
\date{}
\begin{document}
\maketitle

\section{矩阵的秩(补完)}

设 $A\in M_{m,n}(K)$,
\[
A=(\alpha_1,\cdots,\alpha_n)=
\begin{pmatrix}\beta_1\\\vdots\\\beta_m\end{pmatrix}.
\]
上一讲已证 $A$ 的行秩 $=$ $A$ 的列秩, 记作 $A$ 的\textbf{秩} $r(A)$. 这里把"为什么"补齐:

\textbf{(1) 列秩在三种初等列变换下不变}(逐条看列向量组):

\begin{enumerate}
\item 对换两列: $\alpha_1,\cdots,\alpha_i,\cdots,\alpha_j,\cdots,\alpha_n\mapsto\alpha_1,\cdots,\alpha_j,\cdots,\alpha_i,\cdots,\alpha_n$;
\item 某列乘非零常数 $k$: $\alpha_1,\cdots,\alpha_i,\cdots\mapsto\alpha_1,\cdots,k\alpha_i,\cdots$;
\item 某列的倍数加到另一列: $\alpha_j\mapsto\alpha_j+k\alpha_i$.
\end{enumerate}

三种情形下新旧列组都\textbf{线性等价}, 故列秩不变(命题: 线性等价 $\Rightarrow$ 秩相等).

\textbf{(2) 列秩在行变换下不变}: 设 $A=(\alpha_1,\cdots,\alpha_n)\xrightarrow{\ \text{行}\ }A'=(\alpha'_1,\cdots,\alpha'_n)$, 则齐次方程组
\[
AX=0\quad\text{与}\quad A'X=0\ \text{同解},
\]
即
\[
k_1\alpha_1+\cdots+k_n\alpha_n=0
\iff
k_1\alpha'_1+\cdots+k_n\alpha'_n=0 .
\]
于是两者的列组有\textbf{完全相同的线性关系}, 特别地极大无关组的列标相同, 列秩不变.

\textbf{(3) 转置}: $A$ 的列秩 $=$ $A^{T}$ 的行秩. 由 (2) 得"行秩在列变换下不变", 由 (1) 得"行秩在行变换下不变".

\textbf{(4) 化为标准形}: 经初等变换 $A$ 可化为
\[
D=\begin{pmatrix}I_r&0\\0&0\end{pmatrix}\ (m=n),\qquad
D=\begin{pmatrix}I_r&0\end{pmatrix}\ (m<n),\qquad
D=\begin{pmatrix}I_r\\0\end{pmatrix}\ (m>n),
\]
其行秩 $=$ 列秩 $=r$. 由 (1)(2)(3), 初等变换不改变行秩与列秩, 故 $A$ 的行秩 $=$ 列秩 $=r$.

\textbf{从阶梯形看}: $A\xrightarrow{\ \text{行}\ }$ 阶梯形矩阵 $J$. 在 $J$ 中, \textbf{非主元行是主元行的线性组合}(例如 $x_1\beta_1+x_2\beta_2=\beta_3$), 而各主元行线性无关, 故
\[
r(A)=\text{主元个数}=\text{非零行数}.
\]

\section{线性方程组}

\begin{equation*}
\begin{cases}
a_{11}x_1+a_{12}x_2+\cdots+a_{1n}x_n=b_1,\\
a_{21}x_1+a_{22}x_2+\cdots+a_{2n}x_n=b_2,\\
\qquad\qquad\cdots\cdots\\
a_{m1}x_1+a_{m2}x_2+\cdots+a_{mn}x_n=b_m.
\end{cases}
\tag{$*$}
\end{equation*}

要回答三个问题:

(1) 方程组是否有解? 如何判断?

(2) 若有解, 如何求解(Gauss 消元法)?

(3) 解与解之间关系如何?

\begin{defn}[解向量与解集]
如果 $x_1=k_1,\cdots,x_n=k_n$ 是 $(*)$ 的解, 则称 $(k_1,\cdots,k_n)\in K^n$ 是 $(*)$ 的\textbf{解向量}; $(*)$ 的全体解向量构成的集合称为\textbf{解集}.
\end{defn}

\textbf{齐次情形}　若 $b_1=\cdots=b_m=0$, 则 $(*)$ 变成\textbf{齐次线性方程组}
\begin{equation*}
\begin{cases}
a_{11}x_1+\cdots+a_{1n}x_n=0,\\
\qquad\qquad\cdots\cdots\\
a_{m1}x_1+\cdots+a_{mn}x_n=0,
\end{cases}
\tag{$**$}
\end{equation*}
其解集记为 $W\subseteq K^n$.

\begin{prop}
$(**)$ 的解集 $W$ 满足

(1) $\mathbf 0\in W$, 且 $W\neq\varnothing$;

(2) 若 $\eta_1,\eta_2\in W$, 则 $\eta_1+\eta_2\in W$(对加法封闭);

(3) $\forall\eta\in W$, $k\in K$, 则 $k\eta\in W$(对数乘封闭).

于是 \textbf{$W$ 是 $K^n$ 的一个子空间}(齐次解空间).
\end{prop}

\begin{defn}[基础解系]
如果 $W$ 中的向量组 $\eta_1,\cdots,\eta_s$ 满足

(1) $\eta_1,\cdots,\eta_s$ 线性无关;

(2) 任意解向量 $\eta\in W$ 都能被 $\eta_1,\cdots,\eta_s$ 线性表示,

则称 $\eta_1,\cdots,\eta_s$ 是 $(**)$ 的一个\textbf{基础解系}.
\end{defn}

基础解系本质上就是\textbf{解空间 $W$ 的一组基}.

\section{两个预备命题}

\begin{thm}[表示唯一]
若 $\alpha_1,\cdots,\alpha_s$ 线性无关, 且 $\beta\in L(\alpha_1,\cdots,\alpha_s)$, 则 $\beta$ 的表示方式\textbf{唯一}: 若
\[
\beta=k_1\alpha_1+\cdots+k_s\alpha_s=\lambda_1\alpha_1+\cdots+\lambda_s\alpha_s,
\]
则 $k_1=\lambda_1,\cdots,k_s=\lambda_s$.
\end{thm}

\begin{proof}
两式相减得
\[
\mathbf 0=(k_1-\lambda_1)\alpha_1+\cdots+(k_s-\lambda_s)\alpha_s ,
\]
而 $\alpha_1,\cdots,\alpha_s$ 线性无关, 故各系数为零.
\end{proof}

\begin{cor}
$\alpha_1,\cdots,\alpha_s$ 线性无关 $\iff$ 零向量的表示唯一, 即
\[
\mathbf 0=k_1\alpha_1+\cdots+k_s\alpha_s
\]
只有零解 $k_1=\cdots=k_s=0$.
\end{cor}

\section{基础解系的构造}

把 $(**)$ 写成向量形式
\[
x_1\alpha_1+\cdots+x_n\alpha_n=\mathbf 0,\qquad
\alpha_j=\begin{pmatrix}a_{1j}\\\vdots\\a_{mj}\end{pmatrix}\in K^m .
\]

若 $\alpha_1,\cdots,\alpha_n$ 线性无关, 则 $(**)$ \textbf{只有零解}(推论 1).

若 $\alpha_1,\cdots,\alpha_n$ 线性相关, 不妨设 $\alpha_1,\cdots,\alpha_r$($r<n$)是\textbf{极大无关组}. 于是后面 $n-r$ 个向量都能唯一地被它们表示:
\[
\alpha_{r+1}=a_{11}\alpha_1+\cdots+a_{1r}\alpha_r,\quad
\cdots,\quad
\alpha_n=a_{n-r,1}\alpha_1+\cdots+a_{n-r,r}\alpha_r .
\]

令
\[
\eta_1=\bigl(-a_{11},\cdots,-a_{1r},\ 1,\ 0,\cdots,0\bigr),
\]
\[
\eta_2=\bigl(-a_{21},\cdots,-a_{2r},\ 0,\ 1,\cdots,0\bigr),\quad\cdots,\quad
\eta_{n-r}=\bigl(-a_{n-r,1},\cdots,-a_{n-r,r},\ 0,\cdots,0,\ 1\bigr),
\]
即"前 $r$ 个分量取负的表示系数, 后面依次放一个 $1$". 则

\textbf{(1) $\eta_1,\cdots,\eta_{n-r}$ 都是解向量}: 直接代入即得(相当于把 $\alpha_{r+i}$ 的表示移项).

\textbf{(2) $\eta_1,\cdots,\eta_{n-r}$ 线性无关}: 设 $x_1\eta_1+\cdots+x_{n-r}\eta_{n-r}=\mathbf 0$, 由于
\[
x_1\eta_1+\cdots+x_{n-r}\eta_{n-r}
=\bigl(*,\cdots,*,x_1,x_2,\cdots,x_{n-r}\bigr)=\mathbf 0,
\]
比较后 $n-r$ 个分量得 $x_1=\cdots=x_{n-r}=0$.

\textbf{(3) 任一解向量都能被它们表示}: 设 $\eta=(k_1,\cdots,k_r,k_{r+1},\cdots,k_n)$ 是解向量, 令
\[
\eta'=\eta-k_{r+1}\eta_1-\cdots-k_n\eta_{n-r},
\]
则 $\eta'$ 仍是解向量(解空间对线性组合封闭), 且 $\eta'$ 的后 $n-r$ 个分量全为 $0$, 即
\[
\eta'=(c_1,\cdots,c_r,0,\cdots,0).
\]
代入 $(**)$ 得 $c_1\alpha_1+\cdots+c_r\alpha_r=\mathbf 0$; 而 $\alpha_1,\cdots,\alpha_r$ 线性无关, 故 $c_1=\cdots=c_r=0$, 即 $\eta'=\mathbf 0$. 于是
\[
\eta=k_{r+1}\eta_1+\cdots+k_n\eta_{n-r}.
\]

\textbf{结论}　$\eta_1,\cdots,\eta_{n-r}$ 是 $(**)$ 的一个基础解系, 从而
\[
\dim W=n-r=n-r(A).
\]
这就是"\textbf{基础解系所含向量个数 $=n-r(A)$}"的来源.

\end{document}
